\chapter{Interfaces e contratos}\label{cap:interfaces}

\section{Interface do usuário}
A interface usa os mesmos resultados do contrato de serviço. Confirmação de
inscrição só aparece depois da gravação; uma tentativa recusada explica sua
causa sem modificar vagas ou registros.

\begin{quadrosrs}
\caption{Estados observáveis da interface}
\begin{tabelasrs}{L{0.22\textwidth}YY}
\cabecalhosrs\textbf{Estado} & \textbf{Informação ao usuário} & \textbf{Ação disponível}\\\midrule
Catálogo disponível & Nome, início e vagas dos eventos elegíveis. & Selecionar evento.\\
Catálogo vazio & Ausência de eventos elegíveis. & Consultar novamente.\\
Inscrição confirmada & Identificador e evento da inscrição. & Consultar suas inscrições.\\
Inscrição já existente & Identificador ativo e informação de duplicidade. & Abrir o registro existente.\\
Operação recusada & Falta de vagas, prazo encerrado ou sessão inválida. & Escolher outro evento ou autenticar-se, conforme a causa.\\
Falha temporária & Operação sem confirmação e opção de tentar novamente. & Repetir a solicitação com os mesmos dados.\\
\end{tabelasrs}
\fonte{Elaboração didática deste guia.}
\end{quadrosrs}

Mensagens são compreensíveis, preservam dados válidos e não expõem pilhas de
execução ou consultas internas. Os controles respeitam \reqref{RNF-003}.
A tela de inscrições pessoais inclui registros cancelados, distinguindo-os dos
ativos; o botão de cancelamento aparece somente quando a operação é elegível.

\section{Integrações externas}
\begin{quadrosrs}
\caption{Produtores externos e responsabilidade do consumidor}
\begin{tabelasrs}{L{0.25\textwidth}YY}
\cabecalhosrs\textbf{Integração} & \textbf{Dados fornecidos} & \textbf{Validação e falha}\\\midrule
INT-001: catálogo & ID, nome, início em UTC, capacidade e publicação. & Rejeitar dados inválidos ou capacidade inferior às inscrições ativas; conservar a última versão válida.\\
INT-002: identidade & Identidade do participante e sessão com expiração. & Negar operações sem sessão válida; indisponibilidade impede novas sessões, preservando as existentes até expirar.\\
\end{tabelasrs}
\fonte{Elaboração didática deste guia.}
\end{quadrosrs}

A integração usa HTTPS e dados em UTF-8. Os campos de evento obedecem ao
\autoref{tab:dados}. O instante de início recebido é convertido para UTC antes
 de qualquer comparação; o fuso de exibição não altera o prazo.

\section{Contrato HTTP de referência}
O contrato abaixo é didático e independente de framework. As representações
usam JSON e o prefixo \texttt{/api/v1}. As operações pessoais recebem a sessão
pelo cabeçalho \texttt{Authorization}; a identidade é obtida pelo servidor.
Os códigos de resposta seguem a semântica HTTP de \citeonline{http9110}.

\begin{quadrosrs}
\caption{Operações e respostas de sucesso}
\begin{tabelasrs}{L{0.34\textwidth}YY}
\cabecalhosrs\textbf{Método e recurso} & \textbf{Entrada e acesso} & \textbf{Resposta}\\\midrule
GET \texttt{/eventos} & Público; sem corpo. & 200: lista de eventos elegíveis.\\
POST \texttt{/inscricoes} & Sessão válida; \texttt{eventoId}. & 201: novo registro; 200: inscrição ativa existente.\\
GET \texttt{/minhas-inscricoes} & Sessão válida; sem titular fornecido pelo cliente. & 200: registros do titular, inclusive cancelados.\\
GET \texttt{/inscricoes/\{id\}} & Sessão válida; ID de inscrição do titular. & 200: registro autorizado.\\
DELETE \texttt{/inscricoes/\{id\}} & Sessão válida; ID de inscrição do titular, no prazo. & 204: cancelamento realizado ou repetido; sem corpo.\\
\end{tabelasrs}
\fonte{Contrato didático deste guia.}
\end{quadrosrs}

Na criação, a resposta 201 inclui \texttt{Location} para o novo recurso. A
repetição que encontra registro ativo usa 200 e conserva o identificador,
mesmo que o evento já tenha começado: consultar um registro existente não cria
uma nova inscrição. Instantes JSON incluem o sufixo \texttt{Z}, indicando UTC.

As consultas de catálogo e de inscrições pessoais retornam arrays JSON; ausência
de resultados é representada por \texttt{[]}. Cada evento contém \texttt{id},
\texttt{nome}, \texttt{inicio} e \texttt{vagas}. A representação de inscrição,
usada na criação e nas consultas, contém \texttt{id}, \texttt{eventoId},
\texttt{situacao}, \texttt{criadaEm} e \texttt{canceladaEm}; este último campo
é sempre presente e vale \texttt{null} enquanto ativa.

\Needspace{25\baselineskip}
\subsection{Exemplo de criação}
\begin{verbatim}
POST /api/v1/inscricoes
Authorization: Bearer sessao-exemplo
Content-Type: application/json

{"eventoId": "evt-001"}

201 Created
Location: /api/v1/inscricoes/ins-001

{
  "id": "ins-001",
  "eventoId": "evt-001",
  "situacao": "ativa",
  "criadaEm": "2026-10-20T14:00:00Z",
  "canceladaEm": null
}
\end{verbatim}

A resposta de duplicidade contém o mesmo registro com o código 200; a interface
informa que a inscrição já existe. O exemplo pressupõe um evento publicado com
vaga e início posterior ao instante de criação.

\subsection{Erros e precedência}
A sessão é validada antes de qualquer leitura pessoal. Consultar ou cancelar
um ID inexistente ou pertencente a outra pessoa produz a mesma resposta 404.
Após autorizar a operação, a criação verifica duplicidade antes das condições
para novo registro; o cancelamento verifica o prazo antes de alterar o estado.

\begin{quadrosrs}
\caption{Respostas de falha e efeito nos dados}
\begin{tabelasrs}{L{0.09\textwidth}L{0.40\textwidth}Y}
\cabecalhosrs\textbf{HTTP} & \textbf{Código do produto} & \textbf{Condição e efeito}\\\midrule
400 & \texttt{JSON\_INVALIDO} & Corpo não é JSON válido; não alterar dados.\\
401 & \texttt{SESSAO\_INVALIDA} & Sessão ausente ou expirada; negar acesso pessoal.\\
404 & \texttt{RECURSO\_NAO\_ENCONTRADO} & Evento desconhecido ou inscrição ausente/não autorizada; não expor titular.\\
409 & \texttt{SEM\_VAGAS} & Capacidade esgotada para nova inscrição.\\
409 & \texttt{PRAZO\_ENCERRADO} & Instante igual ou posterior ao início.\\
409 & \texttt{EVENTO\_NAO\_PUBLICADO} & Evento não admite nova inscrição.\\
422 & \texttt{DADOS\_INVALIDOS} & Tipo, campo obrigatório ou limite do dicionário não atendido.\\
503 & \texttt{SERVICO\_INDISPONIVEL} & Leitura ou gravação não concluída; não confirmar sucesso.\\
\end{tabelasrs}
\fonte{Contrato didático deste guia.}
\end{quadrosrs}

Erros usam \texttt{codigo} e \texttt{mensagem} no corpo. A resposta 401 também
inclui o cabeçalho \texttt{WWW-Authenticate} com valor \texttt{Bearer}.
O corpo da criação aceita somente \texttt{eventoId}; campos extras são
rejeitados com 422, incluindo um titular indicado pelo cliente.
Uma recusa não altera inscrições nem
vagas. Se a conexão falhar depois da gravação, o cliente pode repetir a criação:
\reqref{RF-003} impede que uma segunda vaga seja consumida.

\section{Hardware e comunicação}
O exemplo não exige sensores ou equipamentos dedicados. Em outro domínio,
registre dispositivos, protocolo, frequência, unidades e resposta à
indisponibilidade. Descreva contratos e resultados necessários, evitando
misturar decisões internas de implementação com obrigações externas do produto.
